\section{Математическая постановка}
\label{sec:problem}

Рассматривается двумерная нестационарная модель теплопереноса в осадочном бассейне, сформулированная в постановке задачи Турнира~\cite{problem2026}. В прямой задаче требуется рассчитать эволюцию поля температуры $T(x,z,t)$ и поля отражательной способности витринита $R_o(x,z,t)$ при заданных начальных и граничных условиях, модели свойств среды и при заданном изменении геометрии осадочного бассейна. Обратная задача заключается в идентификации базального теплового потока $q_{\mathrm{basal}}(x,t)$, приложенного к нижней границе области, по набору наблюдений температуры $T$ и отражательной способности витринита $R_o$ в современный момент времени.

\subsection{Геометрия осадочного бассейна}

В каждый момент времени рассматриваемый двумерный разрез занимает область $\Omega(t)=\{(x,z)|x \in (0, W),\ z \in (z_{\mathrm{surf}}(x,t), L(x,t))\}$, где $W$~--- ширина разреза, $z_{\mathrm{surf}}$ и $L$~--- положения поверхности и подошвы. Координаты $x,z$ выражаются в метрах, ось $z$ направлена вниз. Модельное время $t$ отсчитывается от момента $t_0$, соответствующего древнейшей конфигурации, до $t_{\mathrm{final}}$, соответствующего современности. Геологический возраст $\tau$ отсчитывается в обратном направлении, и современному состоянию соответствует $\tau=0$. Геометрия задана последовательностью конфигураций в моменты $t_k$, а промежуток времени $[t_k,t_{k+1}]$ между двумя соседними конфигурациями называется эпохой. В каждой конфигурации границы слоёв представлены ломаными линиями, и в течение эпохи в процессе осадконакопления возникают новые слои.

В рамках Турнира были предоставлены два набора данных для решения прямой задачи, а именно data1 и data2, называемые далее Моделью~1 и Моделью~2 соответственно, а также набор данных для решения обратной задачи, называемый Моделью~3 (рис.~\ref{fig:geometry}). В наборе данных для обратной задачи отсутствует информация о базальном тепловом потоке $q_{\mathrm{basal}}(x,t)$, поскольку он является искомым параметром, и дополнительно имеются замеры температуры $T$ и отражательной способности витринита $R_o$ на скважинах в современный момент времени.

\reportfigure{01_geometry.png} {Набор данных для решения обратной задачи~--- геометрия осадочного бассейна, литотипы и скважины в современный момент времени. (а) конфигурация для $\tau=200$~млн~лет; (б) конфигурация для $\tau=100$~млн~лет; (в) современная конфигурация для $\tau=0$~млн~лет, расположение скважин и точек измерения.} {fig:geometry}

Внутри эпохи положения кровли и подошвы существующих слоёв линейно интерполируются по времени. Обозначим относительное время эпохи $\theta \in [0, 1]$. Материальная точка нового слоя характеризуется относительной глубиной $a$ в конечной конфигурации, $a=0$ на кровле и $a=1$ на подошве. Положение материальной точки после отложения и момент отложения задаются соотношениями
\begin{gather}
 \theta = \frac{t - t_k}{t_{k+1} - t_k},\qquad
 h(x,t) = \theta h_{\mathrm{end}}(x),\\
 z(a,t) = z_{\mathrm{surf}}(x,t) + h_{\mathrm{end}}(x)\max(0,\theta + a - 1),\\
 t_{\mathrm{dep}} = t_k + (1 - a)(t_{k+1} - t_k).
 \label{eq:deposition}
\end{gather}
Если в следующей конфигурации присутствует новый слой, то процесс осадконакопления моделируется как линейное увеличение мощности слоя с течением времени в предшествующей эпохи, где мощность возрастает от 0 до заданной $h_{\mathrm{end}}(x)$.

При $\theta < 1 - a$~материал считается ещё не отложившимся, и по формуле~\eqref{eq:deposition} занимает положение $z_{\mathrm{surf}}$ на поверхности до момента отложения. Подошва нового слоя откладывается при $t=t_k$, а кровля при $t=t_{k+1}$. Таким образом, моменты отложения различаются по глубине в пределах одного слоя, благодаря чему процесс осадконакопления учитывается точнее, чем если моделировать мгновенное отложение всего слоя целиком. Скорость движения материала $\boldsymbol v$ определяется производной по времени от его положения внутри эпохи. Разломные смещения задаются изменением конфигураций на границах эпох. Изменение геометрии считается известным, и геометрия осадочного бассейна не изменяется при решении обратной задачи.

\subsection{Теплофизические свойства}

На каждой эпохе рассматриваемая область $\Omega(t)$ разбита на несколько подобластей, каждой из которых присвоен свой литотип. Каждая подобласть представляет собой однокомпонентную или двухкомпонентную среду. Для однокомпонентных сред значения теплофизических свойств задаются как функции температуры. Двухкомпонентные среды состоят из скелета породы и порового флюида. Свойства твёрдого скелета которого задаются плотностью $\rho_s$, удельной теплоёмкостью $c_s$, теплопроводностью $\lambda_s$ и объёмной плотностью радиогенного тепловыделения $A_s$. Индекс $f$ обозначает свойства заполняющего поры флюида, а пористость обозначена через $\varphi$. Эффективные теплопроводность $\lambda$, объёмная теплоёмкость $\rho c$ и объёмная плотность радиогенного тепловыделения $A$ определяются следующими зависимостями:
\begin{gather}
 \lambda = \lambda_s(T)^{1 - \varphi}\lambda_f^{\varphi},\qquad
 \rho c = (1 - \varphi)\rho_s c_s(T) + \varphi\rho_f c_f,\qquad
 A = (1 - \varphi)A_s,\\
 c_s(T) = c_{20}\left(0{,}953 + 2{,}29 \cdot 10^{-3} T_C
              - 2{,}835 \cdot 10^{-6} T_C^2 + 1{,}191 \cdot 10^{-9} T_C^3\right),\\
 \lambda_s(T) = 1{,}84 + 358(1{,}0227\lambda_{20} - 1{,}882)
                            \left(\frac{1}{T} - 0{,}00068\right).
\end{gather}
Величина $A$ характеризует мощность тепловыделения при радиоактивном распаде в единице объёма среды и измеряется в Вт/м$^3$. Радиогенные источники тепла сосредоточены в скелете породы, поэтому в выражении для $A$ учитывается его объёмная доля $1-\varphi$.

Здесь температура $T$ выражается в кельвинах, $T_C=T-273{,}15$~--- в градусах Цельсия, а $c_{20}$ и $\lambda_{20}$ обозначают заданные опорные свойства скелета при $20\,{}^\circ$C.

Зависимость теплоёмкости скелета от температуры принята по модели Waples~\cite{waples2004}, а теплопроводности~--- по модели Sekiguchi~\cite{sekiguchi1984}. Изменение пористости $\varphi$ учитывается как по глубине внутри каждого слоя, так и по времени. В каждом слое вводится относительная координата $\alpha=(z-z_{\mathrm{top}})/(z_{\mathrm{bot}}-z_{\mathrm{top}})$, через которую пористость выражается в виде
\begin{equation}
 \varphi(x,z,t) = (1 - \alpha)\varphi_{\mathrm{top}}(x,t)
                  + \alpha\varphi_{\mathrm{bot}}(x,t).
\end{equation}
Значения $\varphi_{\mathrm{top}}$ и $\varphi_{\mathrm{bot}}$ линейно интерполируются между конфигурациями, поэтому свойства материала меняются вслед за его положением, пористостью и температурой.

\subsection{Теплоперенос начальное и граничные условия}

Температура удовлетворяет нестационарному уравнению теплопереноса
\begin{equation}
 (\rho c)\left(\frac{\partial T}{\partial t} + \boldsymbol v \cdot \nabla T\right)
   = \nabla \cdot (\lambda\nabla T) + A,
 \label{eq:heat}
\end{equation}
\begin{gather}
 T\big|_{z = z_{\mathrm{surf}}} = T_{\mathrm{surf}},\qquad
 \lambda\frac{\partial T}{\partial n}\bigg|_{z = L} = q_{\mathrm{basal}},\qquad
 \frac{\partial T}{\partial x}\bigg|_{x = 0,W} = 0,\\
 [T] = 0,\qquad [\lambda\nabla T \cdot \boldsymbol n] = 0.
\end{gather}
На поверхности задаётся температура $T_{\mathrm{surf}}(x,t)$, на фундаменте~--- базальный тепловой поток $q_{\mathrm{basal}}(x,t)$, боковые границы теплоизолированы. На контактах слоёв выполняются условия непрерывности температуры и нормального теплового потока. Здесь $\boldsymbol n$~--- внешняя единичная нормаль к границе, квадратные скобки обозначают скачок величины на внутреннем контакте.

Начальное поле температуры $T_0(x,z)$ определяется как решение стационарного варианта задачи \eqref{eq:heat} в геометрии первой конфигурации
\begin{equation}
 -\nabla \cdot \bigl(\lambda(x,z,0,T_0)\nabla T_0\bigr) = A(x,z,0),
 \qquad T(x,z,0) = T_0(x,z).
\end{equation}
Для расчёта теплофизических свойств используется пористость в начальной конфигурации, зависимость теплофизических свойств от температуры сохраняется, а в качестве граничных условий выступают те же граничные условия, что и в нестационарной задаче~\eqref{eq:heat}, взятые в начальный момент времени. Таким образом, начальное условие согласовано с правой частью и граничными условиями нестационарного уравнения теплопереноса.


\subsection{Отражательная способность витринита}

Зрелость витринита характеризуется показателем отражательной способности $R_o$, выраженным в процентах. В модели EASY\%Ro~\cite{sweeney1990} он определяется совокупностью кинетических реакций, которые протекают независимо друг от друга, в виде
\begin{equation}
  \begin{aligned}
    &I_i(t) = \int_{t_{\mathrm{dep}}}^{t}
    A_i \exp\left(-\frac{E_i}{R_g T(\boldsymbol\xi(s), s)}\right)\,ds,
    \\
    &F = \sum_{i=1}^{N_R} f_i(1 - e^{-I_i}),\qquad
    R_o = \exp(-1{,}6 + 3{,}7F)\quad[\%].
  \end{aligned}
\end{equation}
Для реакции $i$ обозначены предэкспоненциальный множитель $A_i$, энергия активации $E_i$, вес $f_i$ и универсальная газовая постоянная $R_g$. Аррениусовский интеграл $I_i$ скорости реакции характеризует её прогресс и накапливается вдоль траектории материальной точки $\boldsymbol\xi(s)$ начиная с её момента отложения $t_{\mathrm{dep}}$. Для материала начальной конфигурации принимается $t_{\mathrm{dep}}=0$. Начальные интегралы равны нулю, поэтому соответствующее начальное значение $R_o$ определяется самой кинетической моделью. Величина $F$ характеризует взвешенную степень преобразования по всем реакциям.

\subsection{Обратная задача}

В обратной задаче неизвестной величиной принимается базальный тепловой поток $q_{\mathrm{basal}}$, удовлетворяющий следующим условиям:
\begin{enumerate}
  \item Поток не зависит от латеральной координаты $x$, т.е. $q_{\mathrm{basal}}(x,t) = q(t)$;
  \item Поток $q(t)$ монотонно убывает с течением модельного времени $t$;
  \item Поток лежит в пределах $q(t) \in [50,100]$~мВт/м$^2$.
\end{enumerate}
Остальные параметры прямой задачи, такие как теплофизические свойства среды, геометрия осадочного бассейна и другие граничные условия считаются заданными.

Наблюдения температуры и отражательной способности витринита относятся к современному состоянию в момент времени $t_{\mathrm{final}}$. Индекс $w$ обозначает скважину, $j$~--- точку измерения в ней, а индексы $\mathrm{calc}$ и $\mathrm{obs}$~--- расчётное и наблюдённое значения. Число скважин равно $N_w$, количества наблюдений температуры и отражательной способности в скважине $w$ обозначены через $N_w^T$ и $N_w^R$. История потока задаётся вектором параметров $\boldsymbol\theta \in \Theta$, где множество $\Theta$ определяется ограничениями на поток. В соответствии с постановкой задачи Турнира~\cite{problem2026} требуется минимизировать взвешенную сумму квадратов невязок
\begin{equation}
 \begin{aligned}
 J(\boldsymbol\theta) = J_T + J_{R_o}
 &= \sum_{w=1}^{N_w} \sum_{j=1}^{N_w^T}
 w_{w,j}^T \left(T_{w,j}^{\mathrm{calc}}(\boldsymbol\theta) - T_{w,j}^{\mathrm{obs}}\right)^2\\
 &\quad + \sum_{w=1}^{N_w} \sum_{j=1}^{N_w^R}
 w_{w,j}^R \left(R_{o,w,j}^{\mathrm{calc}}(\boldsymbol\theta) - R_{o,w,j}^{\mathrm{obs}}\right)^2.
 \end{aligned}
 \label{eq:objective}
\end{equation}
Здесь $w_{w,j}^T$ и $w_{w,j}^R$~--- положительные веса наблюдений, а $J_T$ и $J_{R_o}$~--- части функционала, отвечающие температуре и отражательной способности. Расчётные значения определяются решением прямой задачи с заданными параметрами потока в точках измерений на конечный момент времени. Искомые параметры определяются условием
\begin{equation}
 \boldsymbol\theta^* = \underset{\boldsymbol\theta \in \Theta}{\operatorname{arg\,min}}\,J(\boldsymbol\theta).
\end{equation}

При известных оценках абсолютных погрешностей измерений $\bar\varepsilon_{w,j}$ и $\bar\eta_{w,j}$ согласование расчёта с данными означает выполнение покомпонентных условий
\begin{gather}
 \left|T_{w,j}^{\mathrm{calc}}(\boldsymbol\theta) - T_{w,j}^{\mathrm{obs}}\right|
 \leq \bar\varepsilon_{w,j},\qquad j = 1,\ldots,N_w^T,\\
 \left|R_{o,w,j}^{\mathrm{calc}}(\boldsymbol\theta) - R_{o,w,j}^{\mathrm{obs}}\right|
 \leq \bar\eta_{w,j},\qquad j = 1,\ldots,N_w^R,
 \qquad w = 1,\ldots,N_w,
 \label{eq:observation-errors}
\end{gather}
что может служить критерием остановки процедуры оптимизации.
